\documentclass[11pt]{article}
\usepackage{fontspec}
\setmainfont{Times New Roman}
\setsansfont{Arial}
\setmonofont{Consolas}
\usepackage{amsmath,amssymb,mathtools}
\usepackage{geometry}
\usepackage{microtype}
\geometry{a4paper,margin=2.35cm}
\setlength{\parindent}{1.25em}
\setlength{\parskip}{0pt}
\makeatletter
\renewcommand\footnotesize{\@setfontsize\footnotesize{7.7}{8.7}\selectfont}
\makeatother
\emergencystretch=35em
\allowdisplaybreaks
\sloppy
\hbadness=10000
\vbadness=10000
\hfuzz=2pt
\vfuzz=10pt
\raggedbottom

\providecommand{\mR}{\mathfrak{R}}
\providecommand{\mA}{\mathfrak{A}}
\providecommand{\mB}{\mathfrak{B}}
\providecommand{\mC}{\mathfrak{C}}
\providecommand{\ma}{\mathfrak{a}}
\providecommand{\mb}{\mathfrak{b}}
\providecommand{\mc}{\mathfrak{c}}
\providecommand{\mm}{\mathfrak{m}}
\providecommand{\mo}{\mathfrak{o}}
\providecommand{\mv}{\mathfrak{v}}

\begin{document}

% Унит: Папер30 Сецтион04 в001. Дирецт Интерславиц Латин транслатион фром the Герман финал-аудитед соурце слице, соурце линес 542--596 / сегментс P30-С0065--P30-С0089. Фоотноте цоунтер цонтинуес афтер Папер30 Сецтион03.
\setcounter{footnote}{14}

\subsection*{\S~4. Теоремы о изоморфији. Директне сумы}

В следујучем предпоставльа се само колцево, односно модулно својство, але никака друга аксиома.

1. Ако $M$ и $\overline M$ сут модулне области взходно к $\mR$ (§2), тогда $M$ называје се хомоморфны к $\overline M$ (точније: модулно-хомоморфны), $M\sim\overline M$, ако всакому елементу из $M$ одповедајучи јест једен и само једен елемент из $\overline M$ тако, же тым $\overline M$ јест вычерпаны;\footnote{Все в смыслу дефиниције равности, ктора влада в $\overline M$; ср. приметку 6.} и ако при тој одповеданьу разлика и множенье тым самым елементом из $\mR$ одповедајучи сут једно другому. То јест, из $\beta\sim\overline\beta$ и $\gamma\sim\overline\gamma$ всегда следује:
\[
(\beta-\gamma)\sim(\overline\beta-\overline\gamma),\qquad r\beta\sim r\overline\beta .
\]

Зато $\mR$-модулу $\mB$ из $M$ хомоморфно одповедајучи јест $\mR$-модул $\overline{\mB}$ из $\overline M$; а целост $\mC^*$ елементов из $M$, кторым одповедајучи сут елементы $\mR$-модула $\overline{\mC}$ из $\overline M$, твори в $M$ модул $\mC^*$, једнозначно определьен чрез $\overline{\mC}$, асоциованы с $\overline{\mC}$, при чем $\mC^*\sim\overline{\mC}$. Ако особливо $\mA$ јест модул, асоциованы с нуловым елементом из $\overline M$, тогда $\mC^*$ јест делитель $\mA$ и разпада се на класы елементов из $M$, конгруентных модуло $\mA$, тако же те класы взаимно једнозначно одповедајучи сут елементам из $\overline{\mC}$. Ако се прејде од $\mB$ в $M$ к $\overline{\mB}$ в $\overline M$, а потом назад к $\mB^*$, добива се $\mB^*=(\mB,\mA)$, то јест највечши обчи делитель $\mB$ и $\mA$; зато $\mB^*=\mB$, ако $\mB$ јест делитель $\mA$.

Ако одповеданье меджу елементами из $M$ и $\overline M$ јест обратно једнозначна, области называјут се изоморфне (точније: модулно-изоморфне), $M\cong\overline M$.

Ако $\mA$ јест либоволны $\mR$-модул из $M$, тогда настава модулна област $\overline M$, хомоморфна к $M$ -- модул класов остатков $M\mid\mA$ -- когда конгруенција по модулу $\mA$ појмаје се како нова релација равности. Всакому елементу из $M$ при том приреджујут се все и само те елементы, кторе сут равне јему в $\overline M$. Преход в $\overline M$ од дефиниције равности к идентитете значи, же все равне елементы в $\overline M$ собирајут се в једну класу -- остаткову класу -- и те остаткове класы појмајут се како нове елементы $\overline M$.

\emph{Всакы хомоморфизм настава преходом к модулу класов остатков;} бо ако $M\sim\overline M$, и ако $\mA$ јест модул асоциованы с нуловым елементом из $\overline M$, тогда, како показано выше, $\overline M$ јест изоморфны к модулу класов остатков $M\mid\mA$.

2. \textbf{Првы теорем о изоморфији.} \emph{Ако $\overline M$ означа модул класов остатков $M\mid\mA$, и ако $\mC$ јест делитель $\mA$, тогда важи изоморфизм
\[
\overline M\mid\overline{\mC}\cong M\mid\mC .
\]}
Бо конгруенција по модулу $\mA$ јест једночасно конгруенција по модулу $\mC$; елементы равне по модулу $\mA$ оставајут равне по модулу $\mC$. Зато можно творити модул класов остатков $M\mid\mC$ -- то јест равност по модулу $\mC$ -- тако, же се најпрво елементы изједначујут по модулу $\mA$, теды преходи се к $\overline M$, и меджу ними собирајут се те, кторе сут равне по модулу $\mC$; в $\overline M$ то одповедајучи јест изједначеньу по модулу $\overline{\mC}$, теды творјеньу $\overline M\mid\overline{\mC}$.

\textbf{Други теорем о изоморфији.} \emph{Ако $\mB$ и $\mA$ сут модулы из $M$, тогда важи изоморфизм
\[
(\mB,\mA)\mid\mA\cong \mB\mid[\mB,\mA].
\]}
Бо по 1. $\mB$ става хомоморфны к $\overline{\mB}$, ако $(\mB,\mA)\mid\mA$ положи се равным $\overline{\mB}$; а понеже при том нуловому елементу в $\overline{\mB}$ одповедајучи сут все елементы из $[\mB,\mA]$ и само они, выше даны изоморфизм зново следује из 1.

3. Ако $\mR$ и $\overline{\mR}$ сут (комутативне) колца, тогда $\mR$ называје се хомоморфно к $\overline{\mR}$ (точније: колцево-хомоморфно), $\mR\sim\overline{\mR}$, ако всакому елементу из $\mR$ одповедајучи јест једен и само једен елемент из $\overline{\mR}$ тако, же тым $\overline{\mR}$ јест вычерпано, и ако при тој приредбе разлике и продукты одповедајучи сут једно другому. Ако одповеданье елементов јест обратно једнозначна, колца называјут се изоморфне (точније: колцево-изоморфне), $\mR\cong\overline{\mR}$.

Все разсматрјаньа из 1. и 2. оставајут валидне, ако модулы заменити идеалами в $\mR$,\footnote{При некомутативных колцах треба взяти за основу двустранне идеалы.} и ако појем модулного хомоморфизма и модулного изоморфизма заменити колцевым хомоморфизмом и колцевым изоморфизмом. Особливо, ако $\mc$ јест делитель $\ma$, колцо класов остатков $\mc\mid\ma$ настава, когда конгруенција по модулу $\ma$ уводи се како нова релација равности; при том треба имети в виду, же тепер сут дефиниране такодже продукты остатковых класов.

Тогда $\mc$ јест хомоморфно к $\mc\mid\ma$; всакы хомоморфизм настава на тој начин, и теоремы о изоморфији важе како в 2:

\textbf{Првы теорем о изоморфији.} Ако $\overline{\mR}$ означа колцо класов остатков $\mR\mid\ma$, и ако $\mc$ јест делитель $\ma$, тогда $\overline{\mR}\mid\overline{\mc}\cong\mR\mid\mc$ в смыслу колцевого изоморфизма.

\textbf{Други теорем о изоморфији.} Ако $\mb$ и $\ma$ сут идеалы из $\mR$, тогда важи колцевы изоморфизм:
\[
(\mb,\ma)\mid\ma\cong\mb\mid[\mb,\ma].
\]\footnote{Те теоремы о изоморфији, знане из теорије груп, најпрво појавјајут се за модулы в троцху специалнејшеј форме у Дедекинда (ср. 3. изд., стр. 484). Выше дана форма за идеалы находи се у Masazo Sono: \emph{Он Цонгруенцес. I, II, III, IV}, Memoirs of the College of Science, Kyoto Империал University, 2 (1917), 3 (1918, 1919), ср. 2, стр. 215. В всаком случају изрично формулованы јест само други теорем о изоморфији.}

4. В следујучем будут собранé рачунске правила за взаимно просте идеалы,\footnote{И то в форме, ктора иде назад к Дедекинду (4. изд., §178, III--VIII). Рачунанье с јединичным елементом, кторо повсуду појавјаје се в новејших изложеньах, много јест сложнејше.} из кторых в 5. будут следовати теоремы о директных сумах. В комутативном колцу $\mR$ предпоставльа се ексистенција јединичного елемента $e$ множеньа. Зато $\ma=\mo\ma$ за всакы идеал из $\mR$, кде под $\mo$ разумеје се јединичны идеал. Два идеалы $\ma,\mb$ называјут се взаимно просте, ако јих највечши обчи делитель $(\ma,\mb)$ става равны $\mo$.

\emph{4$\alpha$. Ако $(\ma,\mb)=\mo$, тогда $(\ma,\mb\mc)=(\ma,\mc)$.} Бо
\[
(\ma,\mc)=(\ma,\mb)\cdot(\ma,\mc)=(\ma^2,\ma\mb,\ma\mc,\mb\mc)
\]
јест делимо чрез $(\ma,\mb\mc)$ и обратно. Из того следује:

\emph{4$\beta$. Из $\mc\mb\equiv0\pmod{\ma}$ и $(\ma,\mb)=\mo$ следује $\mc\equiv0\pmod{\ma}$.} Бо $\mc\mb\equiv0\pmod{\ma}$ јест равнозначно с $(\ma,\mb\mc)=\ma$; зато заради $(\ma,\mb)=\mo$ такодже $(\ma,\mc)=\ma$, или $\mc\equiv0\pmod{\ma}$. Далеј следује:

\emph{4$\gamma$. Ако кажды из идеалов $\ma_1,\ldots,\ma_r$ јест взаимно просты с каждым из идеалов $\mb_1,\ldots,\mb_s$, тогда продукт $\ma_i$ јест взаимно просты с продуктом $\mb_i$.} Бо из $(\ma,\mb)=\mo$ и $(\ma,\mc)=\mo$ добива се $(\ma,\mb\mc)=\mo$, а конечно повторјенье даје тврдженье.

Из 4$\alpha$ и 4$\gamma$ добива се:

\emph{4$\delta$. Ако идеалы $\mb_1,\ldots,\mb_r$ сут попарно взаимно просте, то јест $(\mb_i,\mb_k)=\mo$ за $i\ne k$, тогда јих најменше обче многократно јест равно јих продукту.} Нај $(\ma,\mb)=\mo$, $\mv=[\ma,\mb]$. Тогда
\[
\mv=\mo\mv=(\ma\mv,\mb\mv)=0(\ma,\mb);
\]
понеже $\ma\mb\equiv0\pmod{\mv}$, одтуд следује $\mv=\ma\mb$, а конечно повторјенье с уżываньем 4$\gamma$ даје тврдженье.

\emph{4$\varepsilon$. Ако идеалы $\mb_1,\ldots,\mb_r$ сут попарно взаимно просте, и ако}
\[
\ma_i=[\mb_1,\ldots,\mb_{i-1},\mb_{i+1},\ldots,\mb_r]
=\mb_1\cdots\mb_{i-1}\mb_{i+1}\cdots\mb_r
\]
\emph{јест положено, тогда $(\ma_1,\ldots,\ma_{i-1},\ma_{i+1},\ldots,\ma_r)=\mb_i$, и зато $(\ma_1,\ldots,\ma_r)=\mo$.} Бо по дистрибутивном закону
\[
(\ma_1,\ma_2)=\mb_3\cdots\mb_r(\mb_2,\mb_1)=\mb_3\cdots\mb_r,
\]
и зато
\[
(\ma_1,\ma_2,\ma_3)=\mb_4\cdots\mb_r(\mb_3,\mb_1\mb_2)=\mb_4\cdots\mb_r,
\]
из чего конечным повторјеньем при входном нумерованьу следује тврдженье.

Идеал $\ma_i$ называје се комплементом $\mb_i$ в представјеньу $\mm=[\mb_1,\ldots,\mb_r]$.

5. В комутативном колцу $\mR$ зново предпоставльа се ексистенција јединичного елемента множеньа. Колцо $\mR$ называје се директна сума идеалов $\ma_1,\ldots,\ma_r$ -- в знаких
\[
\mR=\ma_1+\ma_2+\cdots+\ma_r,
\]
ако всакы елемент $c$ из $\mR$ можно представити једным и само једным начином в форме $c=a_1+a_2+\cdots+a_r$, кде всакы $a_i$ јест елемент из $\ma_i$.

\emph{Ако нуловы идеал јест најменше обче многократно попарно взаимно простых идеалов $\mb_1,\ldots,\mb_r$, и ако $\ma_i$ јест комплемент $\mb_i$, тогда $\mR$ јест директна сума идеалов $\ma_1,\ldots,\ma_r$.} Бо заради $\mo=(\ma_1,\ldots,\ma_r)$ всакы елемент из $\mR$ допустја најмене једно адитивно представјенье чрез $\ma_i$; заради $[\ma_i,\mb_i]=(0)$ то представјенье јест једнозначно, бо из $0=a_1+a_2+\cdots+a_r=a_i+b_i$ добива се $a_i=0$. По другом теорему о изоморфији идеалы $\ma_i$ сут изоморфне к колцам класов остатков $\mR\mid\mb_i$. Важе "ортогоналне релације" $\ma_i\ma_k=(0)$ за $i\ne k$ и $\ma_i^2=\ma_i$; последнье заради $\ma_i=\mo\ma_i$.

\emph{Обратно, ако ексистује представјенье $\mR$ како директној сумы,
\[
\mR=\ma_1+\ma_2+\cdots+\ma_r,
\]
и ако положи се
\[
\mb_i=(\ma_1,\ldots,\ma_{i-1},\ma_{i+1},\ldots,\ma_r)
=\ma_1+\cdots+\ma_{i-1}+\ma_{i+1}+\cdots+\ma_r,
\]
тогда $\mb_i$ сут попарно взаимно просте и јих најменше обче многократно јест равно нуловому идеалу.} Бо из $\mo=(\ma_i,\mb_i)$ следује $\mo=(\mb_k,\mb_i)$ за $k\ne i$, понеже $\mb_k$ јест делитель $\ma_i$. Нај далье $c$ јест делимо чрез все $\mb_i$. Тогда
\[
 c=a_2^{(1)}+\cdots+a_r^{(1)}
 =a_1^{(2)}+a_3^{(2)}+\cdots+a_r^{(2)}
 =\cdots=a_1^{(r)}+\cdots+a_{r-1}^{(r)},
\]
кде всакократ $a_i^{(\lambda)}\equiv0\pmod{\ma_i}$. Из једнозначности представјеньа, понеже всакократ једна компонента става нулова, следује $0=a_1^{(\lambda)},\ldots,0=a_r^{(\lambda)}$ за всако $\lambda$, и тым $c=0$.\footnote{Теоремы дане под 5. сут специалне случаји обчего теорема о вези меджу директној сумоју и директным пресеком. Ср. Х. Prüfer, Theorie der Абелсцхен Группен I, Math. Zeitschr. 20 (1924), стр. 165--187, §6. Там дане разсматрјаньа оставајут валидне такодже при тако обчем формулованьу појма групы, же модулы и идеалы подреджујут се јему како специалне случаји.}

\end{document}

